\chapter{Resposta a colisões}

A resposta a colisões é uma etapa fundamental em qualquer sistema de simulação física interativa. Após detectar a interpenetração entre dois corpos, torna-se necessário aplicar um conjunto de correções que restaurem a plausibilidade física do movimento, evitando instabilidades numéricas.

Este capítulo discute os princípios clássicos e as formulações modernas de resposta a colisão, estabelecendo a relação entre os métodos geométricos de detecção e a abordagem baseada em partículas e restrições proposta por \citeonline{jakobsen2001advanced}, que fundamenta este trabalho.

\section{Métodos Dinâmicos na Simulação Física}

A resposta a colisões consiste, essencialmente, em duas operações principais:
\begin{enumerate}
	\item Correção de Posição: eliminar a interpenetração entre dois corpos.
	\item Correção de Velocidade: remover ou ajustar componentes da velocidade que induziriam novo contato imediato.
\end{enumerate}
A literatura apresenta diversas abordagens para modelar o movimento de corpos rígidos e deformáveis. Embora este trabalho se baseie no método simplificado de \citeonline{jakobsen2001advanced}, é importante contextualizar outras categorias amplamente utilizadas em motores físicos modernos:

\subsection*{Método do Impulso}

O Método do Impulso trata colisões aplicando impulsos instantâneos que alteram diretamente as velocidades dos corpos para preservar o momento linear e angular. Essa abordagem é utilizada em motores como Havok~\cite{havok_physics} e Bullet~\cite{bullet_physics}. O impulso é calculado em função da velocidade relativa no ponto de contato, resultando em um método eficiente para simulações em tempo real.

\subsection*{Método de Penalidades}

Nos Métodos de Penalidades, colisões são tratadas como interpenetrações que geram forças de repulsão proporcionais à profundidade de penetração alterando diretamente a aceleração. Essas forças geralmente seguem modelos de mola e amortecimento. Trata-se de um método simples, porém sensível à escolha dos parâmetros de rigidez, podendo causar instabilidade numérica.

\subsection*{Método Baseado em Restrições com Multiplicadores de Lagrange}

Os métodos baseados em restrições formulam os contatos como equações que devem ser satisfeitas exatamente. São resolvidos usando multiplicadores de Lagrange, sendo uma abordagem robusta e adequada para sistemas complexos, mas que exige a solução de sistemas lineares, tornando sua aplicação onerosa em plataformas Web.

\section{Pontos de Contato e Manifold de Colisão}

Em sistemas de simulação física e detecção de colisões, não é suficiente apenas determinar se dois objetos estão  se intersectando. Para a correta resposta física do sistema, é necessário identificar \textit{onde} e \textit{como} ocorre essa interação. Nesse contexto, surgem dois conceitos fundamentais: os pontos de contato (\textit{contact points}) e o manifold de colisão (\textit{collision manifold}).

\subsection*{Pontos de Contato}

Representam as posições no espaço onde dois objetos em colisão efetivamente se tocam. Em termos geométricos, esses pontos correspondem às soluções do problema de proximidade entre superfícies, podendo ser obtidos tanto em situações de interseção quanto em cenários de mínima distância entre objetos convexos.

Em casos de interpenetração, os pontos de contato geralmente são estimados a partir da região de interseção das superfícies, podendo envolver técnicas adicionais, como recorte de faces ou projeções geométricas. Tipicamente, a geometria no instante de contato não é conhecida, mas tem que ser inferida a partir da geometria no instante onde se detecta a interpenetração. Dependendo da configuração dos objetos, pode-se inferir um único ponto de contato (por exemplo, vértice contra face) ou múltiplos pontos (como no contato entre duas faces planas).

\subsection*{Manifold de Colisão}

O manifold de colisão é uma estrutura que agrega e organiza os pontos de contato entre dois objetos em colisão. Em vez de tratar cada ponto de contato de forma isolada, o manifold fornece uma representação coerente da região de contato, incluindo informações adicionais relevantes para a simulação física.

Um manifold de colisão pode ser definido como o conjunto de pontos de contato juntamente com atributos como:
\begin{itemize}
    \item a normal de colisão, que indica a direção de separação entre os objetos;
    \item a profundidade de penetração associada a cada ponto;
    \item possíveis informações auxiliares, como identificadores de arestas ou faces.
\end{itemize}

\section{SAT com Cálculo do MTV}

O SAT pode ser estendido para calcular o vetor de translação mínima (MTV, do inglês \textit{Minimum Translation Vector}). Em vez de encerrar antecipadamente quando uma sobreposição é detectada em um eixo, todos os eixos são testados quanto à sobreposição. Após todos os eixos terem sido testados, o eixo com a menor sobreposição (normalizada) pode ser usado como a normal de contato, e a magnitude da sobreposição pode ser utilizada para estimar a profundidade de penetração
% Com algum trabalho adicional, pontos de contato também podem ser calculados com a ajuda do eixo separador
~\cite{ericson2004real}. Como mostrado no Algoritmo~\ref{alg:sat_mtv}.

% Caso todos os eixos apresentem sobreposição, a menor sobreposição encontrada corresponde à magnitude do vetor necessário para resolver a colisão, como mostrado no Algoritmo~\ref{alg:sat_mtv}, onde $\epsilon$ é um valor numérico para tratar erros de operação flutuante e deve-se ser um valor baixo. 

\begin{algorithm}[htb]
	\caption{SAT com cálculo do MTV}
	\KwIn{Polígonos convexos $A$ e $B$}
	\KwOut{$\vec d$ e $\delta$, ou falso}
	$\vec d \leftarrow \vec 0$ \\
	$\delta \leftarrow \infty$ \\
	\ForEach{aresta $\hat{n}$ de A e B}{
		$\hat{n} \leftarrow$ normal unitária de $\hat{n}$ \\
		$p_1 \leftarrow$ Projeção de $A$ \\
		$p_2 \leftarrow$ Projeção de $B$ \\
		$\Delta \leftarrow$ sobreposição entre $p_1$ e $p_2$ \\
		\If{$\Delta \leq \epsilon$}{
			\Return{False}
		}
		\If{$\Delta < \delta$}{
			$\delta \leftarrow \Delta$ \\
			$\vec d \leftarrow \hat{n}$
		}
	}
	\Return{$\vec d$, $\delta$}
	\label{alg:sat_mtv}
\end{algorithm}

\section{Algoritmo de Expansão de Politopos (EPA)}

O algoritmo de Expansão de Politopos (EPA, do inglês \textit{Expanding Polytope Algorithm}) surge no contexto da evolução dos métodos de detecção de colisão entre objetos convexos, particularmente como uma extensão do algoritmo de GJK. No entanto, sua formulação não fornece diretamente informações necessárias para a resposta à colisão, como a profundidade de penetração ou a direção de separação. O EPA resolve o problema dual do GJK de encontrar a menor translação necessária para separar os objetos, o MTV. O desenvolvimento do EPA é geralmente atribuído a \citeonline{bergen2001proximity}, cujos trabalhos estabeleceram bases importantes para a manipulação geométrica do espaço de configuração (CSO). 

Vimos no GJK que a distância de um par de objetos é um ponto na fronteira do CSO mais próximo da origem. A estratégia para encontrar o MTV é começar com um politopo que contém a origem e tem vértices que se encontram na fronteira do CSO, e expandir esse politopo adicionando vértices que também estão na fronteira. Para isso o GJK pode ser modificado para retornar o simplex final $Q$ e usá-lo como entrada para o EPA. 

O EPA é um algoritmo iterativo em que, para cada etapa da iteração, selecionamos a face do politopo mais próxima da origem. Sejam $p$ o ponto mais próximo da origem que pertence à face selecionada, e $\vec p$ o seu vetor posição. Então, o comprimento de $\vec p$ é um limite inferior para a magnitude da profundidade de penetração, e a normal da face mais próxima da origem é uma aproximação para a direção de separação. Então a face selecionada é subdividida no ponto $p$ e o estende até o ponto de suporte $\vec w = S_{CSO} (\vec p)$. Observe que $||\vec w||$ é um limite \textit{superior} para a profundidade de penetração. O processo continua até que o erro entre $\vec p$ e $\vec w$ seja menor que um valor de tolerância $\epsilon$ ou até que tenham sido feitas aproximações sucessivas suficientes. Veja a Figura~\ref{fig:epa} para uma representação visual do processo de expansão de politopos.

\begin{figure}[htb]
	\centering
	\includegraphics{epa.jpg}
	\caption{Uma sequência de iterações do algoritmo de expansão de politopos. Uma seta indica o vetor mais curto $\vec p$. Um ponto aberto indica o novo vértice $\vec w$. Fonte: Adaptado de~\cite{bergen2001proximity}}
	\label{fig:epa}
\end{figure}

O Algoritmo~\ref{alg:epa} expande iterativamente o simplex final do GJK até encontrar a face da fronteira do CSO mais próxima da origem. A distância $\delta$ da origem a essa face define a profundidade de penetração, enquanto a normal $\hat{n}$ à face define a direção de separação. A tolerância $\epsilon$ (ex: $10^{-4}$) e o limite de iterações são cruciais para evitar loops infinitos em casos de precisão numérica flutuante ou formas curvas aproximadas.

\begin{algorithm}[htb]
    \caption{EPA}
    \label{alg:epa}
    \KwIn{Simplex $Q$ (retornado pelo GJK em caso de colisão)}
    \KwOut{$\hat{n}$, $\delta$}

    \While{iterações $<$ MAX\_ITER}{
        $e = (\vec v_i, \vec v_j) \leftarrow$ aresta mais próxima da origem \\
        $\hat{n}_e \leftarrow$ normal unitária externa de $e$ \\
        $w \leftarrow S_{CSO}(\hat{n}_e)$ \\
        $\delta \leftarrow \vec w \cdot \hat{n}_e$ \\
        \If{$|\delta - d(e)| < \epsilon$}{
            \KwRet{$\hat{n}_e, \delta$}
        }
        Remover a aresta $e$ de $Q$ \\
        Inserir as arestas $(\vec v_i, \vec w)$ e $(\vec w, \vec v_j)$ em $Q$ \\
    }

	\KwRet{$\hat{n}_e, \delta$}
\end{algorithm}

\section{Cálculo dos Pontos de Contato}

Em muitos esquemas de animação física, a detecção de colisão e o cálculo do MTV, utilizando o SAT ou o EPA, são suficientes para realizar a separação dos corpos. Porém, como desejamos manter a coerência geométrica, é importante que haja conservação do momento angular na separação, ou seja, que após a separação os objetos rotacionem, caso necessário. No método de Jakobsen isso é equivalente a mover os vértices dos objetos envolvidos na colisão de forma apropriada. 

Teoricamente, o GJK pode resolver um problema de distância e retornar os pontos de contato diretamente, acelerando essa etapa. Neste trabalho, como nossa implementação distingue os modos de execução entre SAT e GJK e também exigiria grandes modificações no algoritmo, usamos um método adaptado de recorte para ambos os modos de detecção de colisão~\cite{collision_manifolds_2017, dyn4j_clipping_2011}. Originalmente, este método consiste em inferir os primeiros pontos de contatos entre dois objetos convexos, progressivamente recortando a face de um objeto com a borda do segundo objeto e retornando todos os pontos de contato, inclusive aqueles que pertencem a uma aresta. Esse método é comumente usado por bibliotecas como Box2D. 

Neste trabalho adaptamos o método de recorte para retornar apenas os vértices envolvidos na colisão, ou seja, os pontos de contato são sempre vértices. A Figura~\ref{fig:collision-manifold-1} ilustra um caso a analisar onde $\hat n$ é a normal da colisão. 

O primeiro passo consiste em identificar as arestas que fazem parte da colisão. Primeiro selecionamos, para cada objeto, o vértice mais distante ao longo da normal da colisão. Depois, verifica-se os dois vértices adjacentes e selecionamos a aresta que faz o maior ângulo com relação à normal da colisão, como ilustrado na Figura~\ref{fig:collision-manifold-1}. O SAT pode ser modificado para retornar essas informações, ou seja, os vértices e arestas envolvidos na colisão, acelerando essa etapa.

Depois, a face que faz o maior ângulo com relação à normal da colisão é escolhida como a face de referência. A face do outro objeto, que é mais paralela à normal da colisão, é selecionada como a face incidente. Essa etapa é ilustrada na Figura~\ref{fig:collision-manifold-2}. 

\begin{figure}[htb]
    \centering
    \begin{minipage}{0.45\textwidth}
        \centering
		\includesvg[width=\textwidth]{collision-manifold-2}
		\caption{Cenário de colisão, arestas envolvidas em destaque. Fonte: Adaptado de~\cite{collision_manifolds_2017}.}
		\label{fig:collision-manifold-1}
    \end{minipage}
    \hfill
    \begin{minipage}{0.45\textwidth}
        \centering
		\includesvg[width=\textwidth]{collision-manifold-3}
		\caption{Aresta de referência em vermelho e aresta incidente em azul. Fonte: Adaptado de~\cite{collision_manifolds_2017}.}
		\label{fig:collision-manifold-2}
    \end{minipage}
    % \caption{Figuras lado a lado}
\end{figure}

Originalmente, o próximo passo seria recortar a aresta incidente utilizando os planos definidos pela aresta de referência num processo similar ao algoritmo de Sutherland-Hodgman~\cite{collision_manifolds_2017}. Porém, uma vez definidas as arestas de referência e incidente, podemos inferir que os vértices da aresta de referência fazem parte do manifold de colisão e, dos dois vértices da aresta incidente, escolhemos aquele cujo produto escalar com $-\hat n$ é máximo. Então, aplicamos a análise de caso a seguir.

Existem dois casos possíveis de retorno neste método adaptado de recorte. O caso de retornar 3 pontos é o mais comum, onde um vértice de um objeto penetra a face do outro objeto. O caso de retornar 4 pontos é um caso raro de aresta-aresta. O processo de separação é discutido a seguir. Os casos de 0, 1 ou 2 pontos de contato não são possíveis pois sempre detectamos a interpenetração entre dois objetos.

\section{Processo de Separação}

A metodologia de \citeonline{jakobsen2001advanced} pode ser compreendida como uma precursora da moderna dinâmica baseada em posição (PBD, do inglês \textit{Position Based Dynamics})~\cite{muller2007position}. Diferentemente das abordagens que atuam sobre a velocidade ou aceleração, \citeonline{jakobsen2001advanced} trata colisões como restrições geométricas adicionais que devem ser satisfeitas durante o processo iterativo de relaxamento da simulação. Assim, objetos penetrando um ao outro são corrigidos exclusivamente por modificações de posição.

A resposta à colisão envolve dois passos principais. O primeiro consiste em separar os elementos geométricos (vértices, arestas ou faces) que se encontram em interpenetração, o que caracteriza um processo estritamente geométrico. O segundo passo corresponde a um processo iterativo de relaxamento, no qual os elementos afetados ajustam suas posições de acordo com as restrições impostas pelo sistema físico. Existem dois casos principais a considerar: colisão vértice-aresta e colisão aresta-aresta. O caso de colisão vértice-vértice não é possível pois sempre detectamos a interpenetração entre dois objetos.

Considere dois corpos rígidos convexos dinâmicos $A$ e $B$ em estado de interpenetração. 
O subsistema de detecção de colisão fornece:

\begin{itemize}
    \item a normal de separação unitária $\vec n$ (orientada de $A$ para $B$),
    \item a profundidade de penetração $d$,
    \item os vértices envolvidos no primeiro contato.
\end{itemize}

A correção geométrica necessária para eliminar a interpenetração é dada por:

$$
\vec c = d \cdot \vec n
$$

Como ambos os corpos são dinâmicos, a correção é distribuída proporcionalmente às massas dos corpos. 
Sejam $m_A$ e $m_B$ as massas, e $w_A = \frac{1}{m_A}$, $w_B = \frac{1}{m_B}$ suas massas inversas. 
A distribuição global da correção é definida por:
$$
\vec c_A = - \frac{w_B}{w_A + w_B} \vec c, \quad
\vec c_B = \frac{w_A}{w_A + w_B} \vec c
$$
De modo que o corpo de menor massa recebe maior fração da correção, e vice-versa.

\subsection*{Caso Vértice--Aresta}

\begin{figure}[htb]
	\centering
	\includesvg[width=0.45\linewidth]{processo-separacao}
	\caption{Caso colisão vértice-aresta. Correção das posição para configuração válida. O vértice mais próximo de $w$ sofre maior deslocamento, preservando a coerência geométrica da aresta. Fonte: Adaptado de~\cite{jakobsen2001advanced}.}
	\label{fig:processo-separacao}
\end{figure}

Considere um vértice $p \in A$ e uma aresta $\overline{rs} \in B$, definida pelos vértices $r$ e $s$. Seja $w$ um ponto na aresta $\overline{rs}$, correspondente à projeção ortogonal de $p$ sobre a aresta:
$$
\alpha = \frac{(p - r) \cdot (s - r)}{\|s - r\|^2}, \quad \alpha \in [0,1]
$$

$$
w = r + \alpha (s - r)
$$
O objetivo é deslocar $p$ e $w$ para o ponto de contato $w'$, tal que:
$$
w' = w + \vec c_B = p + \vec c_A
$$
A correção aplicada ao vértice $p$ é:
$$
p' = p + \vec c_A
$$

Para a aresta, a correção $\vec c_B$ deve ser distribuída entre os vértices $r$ e $s$ da aresta de forma ponderada pelo parâmetro $\alpha$, de modo que o vértice mais próximo de $w$ receba maior fração do deslocamento.
$$
r' = r + (1 - \alpha)\vec c_B, \quad
s' = s + \alpha \vec c_B
$$

\subsection*{Caso Aresta--Aresta}

Este é um caso raro. Caracterizado por duas arestas $\overline{pq} \in A$ e $\overline{rs} \in B$ que são paralelas. Neste trabalho, optamos por tratar esse caso como uma colisão vértice-aresta selecionando o vértice da aresta incidente mais próximo da aresta de referência para aplicar a correção. 
